\section{Experiments}
\label{sec:experiments}

\subsection{Protocol}

C-MAPSS FD001--FD004 are used with the preprocessing of
Section~\ref{sec:phenomenon} (cap 125, min--max by training statistics,
85/15 train/val window split, last-window test). Metrics are test RMSE
(primary), the asymmetric PHM08 score \cite{saxena2008damage}, $R^2$, and
the failure rate over seeds as flagged by Section~\ref{sec:method_audit}.
Every configuration from Rules 1--4 is trained with five seeds (42, 7, 123,
2026, 3407) controlling initialisation and data order; FD004 additionally
pools four archived runs of the same configuration. We report mean $\pm$ std
over working seeds alongside the failure rate; no C-MAPSS comparison table
we found reports either quantity. Training used an Apple M-series GPU (MPS)
and PyTorch 2.8, at 60--350\,min per run; the 20-run audit took about two
days on one consumer device, and instrumented runs add 50--100\% overhead.
Full settings, the seed-by-seed ledger and the reproducibility file map are
in Supplementary Sections~S1 and~S7.

\subsection{Do the rules hold across seeds}

\begin{table}[t]
\centering\footnotesize
\caption{How often runs failed: (i) the archived runs across all model
variants and (ii) the Rules 1--4 configuration with five seeds per subset.}
\label{tab:audit}
\begin{tabularx}{\linewidth}{Xcccc}
\toprule
Configuration family & FD001 & FD002 & FD003 & FD004 \\
\midrule
Archive, all labelled variants (157 runs) & 0/45 & 16/43 & 0/36 & 8/33 \\
Rules 1--4, five seeds & 0/5 & 0/5 & 0/5 & 6/8 \\
\bottomrule
\end{tabularx}
\end{table}

On FD001, FD002 and FD003 the chosen settings worked in all fifteen runs,
with per-seed RMSE tight enough for mean $\pm$ std to be meaningful:
13.66 $\pm$ 0.10, 19.43 $\pm$ 0.44 and 16.09 $\pm$ 0.52
(Table~\ref{tab:audit}, Fig.~\ref{fig:seeds}). On FD004 the seed decides: six
of eight runs ended at the fixed output predicted by
Eq.~(\ref{eq:fingerprint}), while the two that worked reached 19.46/19.60
(std 0.07), with nothing in between. The rules remove the settings that
always failed (FD002: 37\% in the archive, 0/5 here) but cannot make every
seed work everywhere; the checks in Section~\ref{sec:method_audit} catch the
remaining failures in minutes and let us start a fresh run. No audited FD004
configuration trained reliably over multiple seeds, so this looks like a
property of that region rather than of one poor rule.

\begin{figure}[t]
\centering
\includegraphics[width=0.56\linewidth]{figures/fig9_seed_boxplot.png}
\caption{Per-seed test RMSE for the Rules 1--4 configuration (teal circles =
working runs; coral crosses = failed runs; navy line = mean of the working
runs). FD004 splits into two groups: six seeds at 43.67--43.72, two at
19.46--19.60, nothing in between.}
\label{fig:seeds}
\end{figure}

\subsection{Accuracy, score, and component removals}

The chosen configuration beats a vanilla Transformer on every subset and
sits in the DCNN/BiGRU-TSAM/MS-DCNN band, but stays below specialised
architectures such as DKAMFormer \cite{kim2025dkam}, which targets
FD002/FD004 with dedicated multi-scale machinery. The full comparison
against twelve published methods is in Supplementary Table~S5; none of
those rows reports a failure rate. Mean PHM08 scores over working seeds are
307 (FD001), 2\,855 (FD002), 1\,050 (FD003) and 1\,783--2\,448 (FD004),
against 61\,569--75\,860 for failed runs. The asymmetric penalty therefore
multiplies the fixed-output failure by 15--250$\times$ relative to RMSE,
which is why Score, not RMSE, is the metric that matters in service.

The eight-way component removal (A0--A7, four subsets, with PHM08 scores
recomputed from saved checkpoints) is in Supplementary Table~S6. Read
alongside Section~\ref{sec:mechanism}, it carries much of the causal
argument: taking out a branch the model was thought to depend on (A1 on
FD002; A3/A4 on FD004) produces the same fixed-output failure confirmed by
Eq.~(\ref{eq:fingerprint}), rather than a graceful loss of accuracy, whereas
an accuracy-only reading would rank A5, where removing max pooling improves
the mean, as the best change. The instrumented runs supply the complement:
recovery is driven by the TCN branch (a hundredfold gradient jump,
$\times$129--235), the same branch whose removal forces FD002 to fail for
good.

Supplementary Sections~S5--S6 also give the segment-error signature on every
subset (FD002: 69.9 / 38.5 / 10.7 / 38.2 across RUL buckets $<30$, 30--60,
60--100, $\geq100$), the val--test gaps of working runs ($-2.3$\% to
$+25.3$\%, largest on FD003, a known limitation of windowed splits retained
for comparability), and the layer-1 attention maps (mean normalised row
entropy 0.7221 for a working checkpoint against exactly 1.0000 for a failed
one).
